\documentclass[10pt,a4paper]{amsart}
\usepackage{amsmath,amssymb,amsthm,graphicx}
\usepackage{verbatim}
\newenvironment{reference}{\comment}{\endcomment}
%\newenvironment{reference}{}{}
\newenvironment{slogan}{\comment}{\endcomment}
\newenvironment{history}{\comment}{\endcomment}

% For commutative diagrams we use Xy-pic
\usepackage[all]{xy}

% We use 2cell for 2-commutative diagrams.
\xyoption{2cell}
% Theorem environments.
%
\theoremstyle{plain}
\newtheorem{theorem}[subsection]{Theorem}
\newtheorem{proposition}[subsection]{Proposition}
\newtheorem{lemma}[subsection]{Lemma}

\theoremstyle{definition}
\newtheorem{definition}[subsection]{Definition}
\newtheorem{example}[subsection]{Example}
\newtheorem{exercise}[subsection]{Exercise}
\newtheorem{situation}[subsection]{Situation}

\theoremstyle{remark}
\newtheorem{remark}[subsection]{Remark}
\newtheorem{remarks}[subsection]{Remarks}

\numberwithin{equation}{subsection}

% Macros
%
\def\lim{\mathop{\mathrm{lim}}\nolimits}
\def\colim{\mathop{\mathrm{colim}}\nolimits}
\def\Spec{\mathop{\mathrm{Spec}}}
\def\Hom{\mathop{\mathrm{Hom}}\nolimits}
\def\Ext{\mathop{\mathrm{Ext}}\nolimits}
\def\SheafHom{\mathop{\mathcal{H}\!\mathit{om}}\nolimits}
\def\SheafExt{\mathop{\mathcal{E}\!\mathit{xt}}\nolimits}
\def\Sch{\mathit{Sch}}
\def\Mor{\mathop{\mathrm{Mor}}\nolimits}
\def\Ob{\mathop{\mathrm{Ob}}\nolimits}
\def\Sh{\mathop{\mathit{Sh}}\nolimits}
\def\NL{\mathop{N\!L}\nolimits}
\def\CH{\mathop{\mathrm{CH}}\nolimits}
\def\proetale{{pro\text{-}\acute{e}tale}}
\def\etale{{\acute{e}tale}}
\def\QCoh{\mathit{QCoh}}
\def\Ker{\mathop{\mathrm{Ker}}}
\def\Im{\mathop{\mathrm{Im}}}
\def\Coker{\mathop{\mathrm{Coker}}}
\def\Coim{\mathop{\mathrm{Coim}}}

% Boxtimes
%
\DeclareMathSymbol{\boxtimes}{\mathbin}{AMSa}{"02}

%
% Macros for moduli stacks/spaces
%
\def\QCohstack{\mathcal{QC}\!\mathit{oh}}
\def\Cohstack{\mathcal{C}\!\mathit{oh}}
\def\Spacesstack{\mathcal{S}\!\mathit{paces}}
\def\Quotfunctor{\mathrm{Quot}}
\def\Hilbfunctor{\mathrm{Hilb}}
\def\Curvesstack{\mathcal{C}\!\mathit{urves}}
\def\Polarizedstack{\mathcal{P}\!\mathit{olarized}}
\def\Complexesstack{\mathcal{C}\!\mathit{omplexes}}
% \Pic is the operator that assigns to X its picard group, usage \Pic(X)
% \Picardstack_{X/B} denotes the Picard stack of X over B
% \Picardfunctor_{X/B} denotes the Picard functor of X over B
\def\Pic{\mathop{\mathrm{Pic}}\nolimits}
\def\Picardstack{\mathcal{P}\!\mathit{ic}}
\def\Picardfunctor{\mathrm{Pic}}
\def\Deformationcategory{\mathcal{D}\!\mathit{ef}}
\usepackage{cleveref}
\crefname{originalreference}{original reference}{original references}
\Crefname{originalreference}{Original reference}{Original references}
\makeatletter
\newlabel{varieties-lemma-dimension-fibres-locally-algebraic}{{0B2L}{source}{}{sourcepointer-1}{}}
\newlabel{varieties-lemma-dimension-fibres-locally-algebraic@cref}{{[originalreference][0][]0B2L}{[1][0][]source}{}{}{}}
\makeatother
\title{General projections: the excess of Y cap r\_p\textasciicircum{}\{-1\}(r\_p(Z)) over Y cap Z is bounded}
\author{The Stacks Project Authors}
\date{stacks-project:intersection.tex:local-snapshot-20261012; GNU FDL 1.2 or later}
\begin{document}
\maketitle
\noindent\textit{Editorial note.} All delivered mathematical text is the author's own, copied verbatim from the pinned Stacks Project chapter file intersection.tex; no statement or proof step was rewritten, reordered or removed.\par\medskip
\noindent\textit{Editorial note.} Span recomputation by label pairing (LANE\_BRIEF 10.2): in this chapter the author's proof environment follows the statement environment as a standalone block; the spec covers the complete environment pair (statement span, proof span) as recorded in the pool candidate.\par\medskip
\noindent\textit{Editorial note.} Occupation scan: combined claim\_key + statement\_anchor + generic-line scan over all registration files (active and withdrawn) found 0 hits for this chapter (is-occupied.json, 2026-10-14); consistent with the duty manager's six-check dispatch.\par\medskip
\ifdefined\hypertarget\hypertarget{sourcepointer-1}{}\fi
\noindent\textit{Source reference.} Original reference 0B2L: sources/stacks-chow/tags.txt, source lines 13934--13934. Cross-chapter standard prerequisite cited by the author inside the counted span; it is used as a pointer and its own proof is not part of this item.\par\medskip
\begin{lemma}
\label{lemma-projection-injective}
Let $V$ be a vector space of dimension $n + 1$.
Let $Y, Z \subset \mathbf{P}(V)$ be closed subvarieties.
There is a nonempty Zariski open $U \subset \mathbf{P}(V)$
such that for all closed points $p \in U$ we have
$$
Y \cap r_p^{-1}(r_p(Z)) = (Y \cap Z) \cup E
$$
with $E \subset Y$ closed and
$\dim(E) \leq \dim(Y) + \dim(Z) + 1 - n$.
\end{lemma}

\begin{proof}
Set $Y' = Y \setminus Y \cap Z$.
Let $y \in Y'$, $z \in Z$ be closed points with $r_p(y) = r_p(z)$.
Then $p$ is on the line $\overline{yz}$ passing through $y$ and $z$.
Consider the finite type scheme
$$
T = \{(y, z, p) \mid y \in Y', z \in Z, p \in \overline{yz}\}
$$
and the morphism $T \to \mathbf{P}(V)$ given by $(y, z, p) \mapsto p$.
Observe that $T$ is irreducible and that $\dim(T) = \dim(Y) + \dim(Z) + 1$.
Hence the general fibre of $T \to \mathbf{P}(V)$ has dimension at most
$\dim(Y) + \dim(Z) + 1 - n$, more precisely, there exists a nonempty
open $U \subset \mathbf{P}(V) \setminus (Y \cup Z)$ over
which the fibre has dimension at most $\dim(Y) + \dim(Z) + 1 - n$
(Varieties, Lemma \ref{varieties-lemma-dimension-fibres-locally-algebraic}).
Let $p \in U$ be a closed point and let $F \subset T$ be the fibre
of $T \to \mathbf{P}(V)$ over $p$. Then
$$
(Y \cap r_p^{-1}(r_p(Z))) \setminus (Y \cap Z)
$$
is the image of $F \to Y$, $(y, z, p) \mapsto y$. Again by
Varieties, Lemma \ref{varieties-lemma-dimension-fibres-locally-algebraic}
the closure of the image of $F \to Y$ has dimension at most
$\dim(Y) + \dim(Z) + 1 - n$.
\end{proof}

\end{document}
